\subsection{Cost Analysis: Real Data Collection vs Synthetic Generation}
\label{subsec:cost-analysis}

A key practical consideration for adopting synthetic data is the cost trade-off between collecting additional real traces versus generating synthetic data. Table~\ref{tab:cost-comparison} compares the computational cost of synthetic trace generation across different context lengths.

\begin{table}[t]
\centering
\caption{Computational cost of synthetic trace generation (GPU hours on H100). Costs scale with context length but remain modest compared to real data collection infrastructure.}
\label{tab:cost-comparison}
\small
\begin{tabular}{lrrr}
\toprule
\textbf{Benchmark} & \textbf{L=256} & \textbf{L=1024} & \textbf{L=4096} \\
\midrule
ffmpeg & 10 & 10 & 4 \\
iozone & 24 & 16 & 16 \\
pybench & 5 & 1 & 0.2 \\
scimark2 & 6 & 1 & 0.2 \\
stream & 13 & 6 & 3 \\
unpack-linux & 24 & 13 & 3.9 \\
\midrule
\textbf{Total (6 benchmarks)} & \textbf{82} & \textbf{47} & \textbf{27.3} \\
\bottomrule
\end{tabular}
\end{table}

\paragraph{Synthetic Generation Cost.} Generating synthetic traces for all six benchmarks at L=4096 (optimal quality) requires approximately 27 GPU hours on an H100. At current cloud pricing (\$2-3/hour for H100 instances), the total cost is \textbf{\$54-82 for a complete synthetic dataset}. This is a one-time cost that produces unlimited synthetic traces for augmentation.

\paragraph{Real Data Collection Cost.} In contrast, collecting real kernel traces requires: (1) instrumentation infrastructure (LTTng setup, storage, monitoring), (2) dedicated hardware for reproducible trace collection, (3) multiple runs per benchmark to ensure statistical validity, and (4) engineering time for setup, execution, and validation. For a production system, this typically involves:
\begin{itemize}[nosep,leftmargin=*]
    \item Infrastructure setup: 40-80 engineering hours (\$4,000-8,000 at \$100/hour)
    \item Compute resources: Dedicated machines for trace collection (\$500-1,000/month)
    \item Storage: 100GB-1TB of trace data (\$50-200/month)
    \item Ongoing maintenance: 10-20 hours/month (\$1,000-2,000/month)
\end{itemize}

\paragraph{Cost-Benefit Analysis.} For organizations needing to double their training dataset, synthetic generation provides a \textbf{50-100× cost reduction} compared to collecting equivalent real data. Even accounting for the initial diffusion model training cost (200-400 GPU hours, \$400-1,200), synthetic generation becomes cost-effective after the first use. For privacy-sensitive environments where real data collection is infeasible, synthetic generation provides the only viable path to dataset expansion.

\paragraph{Scalability.} Synthetic generation costs scale sub-linearly with dataset size: generating 10× more synthetic traces requires minimal additional cost (just inference time), whereas collecting 10× more real data requires proportional infrastructure and engineering effort. This makes synthetic data particularly attractive for large-scale deployments or scenarios requiring frequent dataset updates.